\section{Exercise 4: Addition}
\subsection{Task and method}
Implement \texttt{[yn, yorigin] = add(x1n, x1origin, x2n, x2origin)}
to obtain
\begin{equation}
  y(n)=x_1(n)+x_2(n).
\end{equation}
The samples being added must refer to the same time $n$. Equal vector
lengths alone do not ensure that the samples are aligned, since the
two input origin positions may differ.

Let $N_1,N_2$ denote the vector lengths and $o_1,o_2$ their origin
parameters. The stored time intervals are
\[
  1-o_1\leq n\leq N_1-o_1,\qquad
  1-o_2\leq n\leq N_2-o_2.
\]
Construct a common interval with endpoints
\begin{align}
  n_{\min}&=\min(1-o_1,1-o_2),\\
  n_{\max}&=\max(N_1-o_1,N_2-o_2).
\end{align}
Initialize two vectors of length $n_{\max}-n_{\min}+1$ with zeros.
Place each input sample at its common-vector position
\begin{equation}
  i=n-n_{\min}+1.
\end{equation}
Samples outside each input's stored interval remain zero. The aligned
vectors can then be added element by element. The output origin is
\begin{equation}
  \texttt{yorigin}=1-n_{\min}.
\end{equation}
The implementation displays both aligned inputs and their sum in three
subplots within the same figure, using shared axis limits.

\subsection{Source code}
\lstinputlisting[caption={Scilab addition function and example call},label={lst:add}]{../4.sci}

\subsection{Example and expected result}
The handout specifies the example
\begin{lstlisting}[numbers=none]
[yn, yorigin] = add([0, 1, 3, -2], 1, [1, 1, 2, 3], 2);
\end{lstlisting}
The first vector covers $n=0,1,2,3$, whereas the second covers
$n=-1,0,1,2$. The common time interval is therefore $-1\leq n\leq3$.
\begin{table}[H]
  \centering
  \caption{Time alignment and samplewise addition}
  \label{tab:add}
  \begin{tabular}{lrrrrr}
    \toprule
    Time $n$ & $-1$ & 0 & 1 & 2 & 3\\
    \midrule
    Aligned $x_1(n)$ & 0 & 0 & 1 & 3 & $-2$\\
    Aligned $x_2(n)$ & 1 & 1 & 2 & 3 & 0\\
    Sum $y(n)$ & 1 & 1 & 3 & 6 & $-2$\\
    \bottomrule
  \end{tabular}
\end{table}
The calculated output agrees with the result stated in the handout:
\[
  \texttt{yn}=[1,1,3,6,-2],\qquad \texttt{yorigin}=2.
\]

\subsection{Execution evidence}
\resultfigure{images/exercise4.png}
  {Aligned input signals and their sum for Exercise 4}
  {fig:add}
The captured Scilab plot agrees with the time-aligned sum in
\Cref{tab:add}, including the output samples at both ends of the interval.

\subsection{Discussion}
At $n=2$, the two samples have amplitudes $3$ and $3$, giving an output
of $6$. At $n=-1$, only the second input contributes, so $y(-1)=1$.
At $n=3$, only the first input contributes, so $y(3)=-2$.
The output contains five stored samples even though each input contains
four. Directly adding the original vectors would incorrectly pair
samples at different time indices and would not represent the required
signal sum.
